% ============================================================================
% CHAPITRE 2 : INTRODUCTION
% ============================================================================

% ----------------------------------------------------------------------------
\section{INTRODUCTION}
\label{sec:introduction_cadre_formel}
% ----------------------------------------------------------------------------

\begin{sloppypar}
La th\`ese progresse selon un cycle canonique en six actes,
\textbf{F}$\to$\textbf{B}$\to$\textbf{E}$\to$\textbf{S}$\to$\textbf{C}$\to$\textbf{X}.
Ce chapitre en r\'ealise les trois premiers. La section~\ref{sec:clo}
formalise le cadre CLO et d\'efinit les trois m\'etriques
($\Conf$, $\Rel$, $\Opp$) qui rendent le paradoxe tractable, en y
int\'egrant les jeux cryptographiques (IND-CONF, BIND-REL, EXT-OPP) qui
captent leur contenu op\'erationnel. La section~\ref{sec:cro_classique}
borne les compromis r\'ealisables dans le r\'egime classique par
le Trilemme CRO conditionnel. La section~\ref{sec:cro_pq} \'etend
ces bornes aux adversaires quantiques et aux canaux post-quantiques. La
section~\ref{sec:conclusion_cadre_formel} synth\'etise ces trois
contributions. Les trois actes suivants, saturer, certifier et
exp\'erimenter, font l'objet des chapitres~\ref{chap:constructions}
et~\ref{chap:experimentation}.

Ce chapitre \'etablit le socle th\'eorique de la th\`ese. Chaque section
correspond \`a une contribution publi\'ee et apporte une r\'eponse pr\'ecise
\`a l'un des objectifs de recherche formul\'es \`a la
sous-section~\ref{subsec:objectifs_recherche}.

% ----------------------------------------------------------------------------
\subsection{D\'emarche m\'ethodologique~: les actes
\textbf{F}, \textbf{B} et \textbf{E}}
\label{subsec:demarche_methodologique}
% ----------------------------------------------------------------------------

Avant d'entrer dans les sections techniques, nous exposons la logique de
chaque \'etape. Dans le cycle
\textbf{F}$\to$\textbf{B}$\to$\textbf{E}$\to$\textbf{S}$\to$\textbf{C}$\to$\textbf{X},
chaque \'etape r\'epond \`a une n\'ecessit\'e m\'ethodologique sp\'ecifique.

\paragraph{Formaliser (\textbf{F}).}
La cryptographie institutionnelle ne disposait d'aucun cadre m\'etrique
pour l'interpr\'etabilit\'e. Avant de borner, il fallait d\'efinir ce
qui doit \^etre born\'e. La m\'ethodologie adopt\'ee est
axiomatique-constructive. Elle part de la th\'eorie de l'information
(Shannon, Fano, Pinsker) pour d\'efinir des m\'etriques mesurables,
comparables et formellement manipulables. La formalisation pr\'ec\`ede
les bornes parce que sans d\'efinition pr\'ecise de l'opposabilit\'e,
toute borne serait circulaire.

\paragraph{Borner (\textbf{B}).}
Une fois le cadre CLO \'etabli, il reste \`a d\'eterminer les limites
fondamentales du r\'ealisable. La m\'ethodologie des
th\'eor\`emes d'impossibilit\'e conditionnelle, inspir\'ee des trilemmes
CAP et FLP, conditionne les bornes aux param\`etres cognitifs r\'eels des
institutions ($\Delta$, $\alpha$, $\beta$). Cette approche produit une
contrainte param\'etr\'ee par les acteurs institutionnels plut\^ot qu'un
obstacle th\'eorique absolu. Lorsque la granularit\'e des
verdicts d\'epasse la capacit\'e d'extraction institutionnelle $\Delta$, c'est
$\Delta$ qui plafonne les compromis r\'ealisables~; sinon, c'est l'in\'egalit\'e
de Fano.

\paragraph{\'Etendre (\textbf{E}).}
L'extension quantique introduit une perturbation qualitativement
diff\'erente, qui tient au bruit de canal quantique et \`a la perte
d'interpr\'etabilit\'e qu'il induit lorsque la preuve est port\'ee par un support quantique. Les outils de la th\'eorie
de l'information quantique, entropie de von Neumann, borne de Holevo (sous-sous-section~\ref{subsubsec:holevo}),
lemme de Gentle Measurement, permettent d'\'etendre les bornes \`a cet
environnement sans repartir de z\'ero. La borne explicite conserve sa forme,
avec une capacit\'e quantique $\Delta_q$~; le coefficient de perte
d'interpr\'etabilit\'e $\eta_q$ n'appara\^it que dans une majoration
affaiblie. Ce coefficient est nul pour les preuves qui sont des cha\^ines de bits.
\end{sloppypar}
